(a) A change of basis of $\mathcal E$ induces an automorphism of its symmetric algebra whose degree-one part is the same invertible linear change. Thus the affine-space charts of $\operatorname{Spec}S(\mathcal E)$ have linear transition maps. The identity on $\operatorname{Spec}S(\mathcal E)$ gives an isomorphism between the bundle structures obtained from different bases, since transitions between any two such charts remain linear.

(b) Over $V=\operatorname{Spec}A$ in a bundle chart, a section corresponds to an $A$-algebra map $A[x_1,\ldots,x_n]\to A$, hence to the tuple of images of the $x_i$. Addition and scalar multiplication of tuples give the sheaf of sections the structure $\mathcal O_V^n$. The transitions are linear, so these module structures agree on overlaps and make the sheaf locally free of rank $n$.

(c) The universal property of the symmetric algebra identifies $\operatorname{Hom}_{\mathcal O_V}(\mathcal E|_V,\mathcal O_V)$ with the $\mathcal O_V$-algebra maps $S(\mathcal E|_V)\to\mathcal O_V$. By \exref{II.5.17}, these correspond to sections $V\to\mathbb V(\mathcal E)|_V$. Locally this is the tuple identification in (b), so it is linear and compatible with restriction. Hence $\mathcal S(\mathbb V(\mathcal E)/Y)\cong\mathcal E^\vee$.

(d) Associate to a bundle its sheaf of sections $\mathcal S$, and then $\mathcal S^\vee$. Part (c) and the double dual isomorphism recover $\mathcal E$ from $\mathbb V(\mathcal E)$. Conversely, the affine-space charts of a bundle identify it with $\mathbb V(\mathcal S^\vee)$: their linear coordinate changes are exactly those induced on the dual of the sheaf of sections. The identifications glue, giving the required bijection on isomorphism classes.
